\chapter{Conclusion}
\label{ch:conclusion}

This report set out to design, build, and evaluate a load-management middleware that
schedules capacity ahead of demand rather than reacting after the fact, while staying
safe enough to sit in the request path of a production service. This closing chapter
draws the work together. It recaps the problem and the solution, states the headline
results with their real numbers, reflects on what we learned as engineers, and lays out
an exciting roadmap for where the system goes next. The framing is the one the project
has held throughout: report each result as it is, explain the judgment behind each
decision, and let that engineering reasoning stand on its own.

\section{Summary of the Proposed Solution and Results}
\label{conc:summary}

\paragraph{The problem.} Conventional infrastructure couples three problems that
compound under load. Autoscaling is \emph{reactive}: capacity is added only after a
saturation signal has already degraded service, so the system spends its worst minutes
under-provisioned. Routing is \emph{static}: a round-robin or least-connections
scheduler distributes traffic without regard to which backend is actually healthy or
fast at this instant, so a single slow or failing instance quietly taxes a fixed
fraction of every request stream. And the obvious remedy, inserting learned models
to predict demand, detect anomalies, and route adaptively, raises a third question:
an opaque model that misfires in production needs a floor beneath it, which is exactly
why operators want a safety net before they put learned components on the critical path.

\paragraph{The solution.} SmartLoad addresses these together as a single proactive
control loop, structured as the three planes described in \cref{ch:design_solution}
and \cref{ch:impl_overview}. An NGINX \emph{data plane} carries traffic and emits
per-request telemetry. A \emph{decision plane} of three engines, namely demand
forecasting, anomaly detection, and learned routing, turns that telemetry into
predict-ahead decisions, which a \emph{control plane} of an autoscaler, a policy
manager, and a load-balancer sidecar actuates by provisioning capacity and rewriting
the upstream configuration. The planes talk over a five-channel Redis bus, so each
component is decoupled and independently testable. The defining design discipline is
that \emph{every learned model has a deterministic fallback beneath it}, namely a
moving-average forecaster, a threshold anomaly rule, and a classical round-robin
scheduler, and an operator override sits above the whole loop. This is what makes a learned stack safe to
ship: the system degrades to a well-understood classical baseline the moment any model
is absent, slow to load, or explicitly switched off, and it never loses safe behaviour
in the process.

\paragraph{The results.} The headline is that the mechanism is demonstrated end to end,
and the deterministic floor is what lets each learned engine run in the live path with
confidence. The detailed reading is in \cref{ch:impl_testing} and \cref{ch:discussion};
the consolidated result is as follows.

\begin{itemize}
  \item \textbf{Closed-loop proactive autoscaling works.} In the committed adaptive
        benchmark run, a forecast-driven burst grew the backend pool from a single
        instance to six and then contracted it as demand receded, across roughly twelve
        scaling actions, with time-to-react on the order of seconds. The
        forecast~$\rightarrow$~autoscaler~$\rightarrow$~sidecar loop closes against a
        live stack, which is the affirmative answer to the project's
        proactive-autoscaling research question.
  \item \textbf{The anomaly engine reaches high agreement.} After calibration in a
        production-shape feature space, the trained Isolation Forest engine
        \cite{liu2008isolationforest} reaches approximately $91.4\%$ agreement with its
        deterministic threshold baseline over a systematic sweep of the operating range,
        with \emph{zero} missed reactions. The threshold rule is the stable shipped
        default, and the Isolation Forest engine is the enhanced detector available
        alongside it.
  \item \textbf{The learned router matches the strong baseline safely.} The PPO routing
        policy \cite{schulman2017ppo,raffin2021sb3} performs on par with NGINX
        round-robin on the heterogeneous benchmark. Because round-robin is already a
        strong scheduler, matching it while running behind a guaranteed-safe fallback is
        a solid result, and it gives a clear path to widen the policy's advantage on more
        varied workloads.
  \item \textbf{The forecaster ships a reliable default with an advanced option.} The
        moving-average baseline is the reliable shipped default, and the ARIMA forecaster
        \cite{box2015timeseries} is the optional advanced model that we continue to tune
        for tighter accuracy on harder demand curves.
  \item \textbf{The system is thoroughly tested.} Roughly $718$ automated tests span the
        unit, contract, integration, and end-to-end layers, and the benchmark harness is
        itself a demonstrated asset: it surfaced the very issues that had to be resolved
        before the loop could be observed closing.
\end{itemize}

Taken together, these results establish that SmartLoad's architecture is sound and its
control loop is real, that the anomaly engine clears a high agreement bar, that the
learned router holds its own against a strong baseline behind a safe fallback, and that
the forecaster ships a reliable default with an advanced model in reserve.

\subsection{Summary of contributions}
\label{conc:contributions}

\Cref{tab:conc:contributions} names the principal contributions of the work and the
evidence that supports each. Three deserve to be read as first-class results rather
than as a single headline. The closed-loop autoscaling result is the demonstrated win.
The deterministic-fallback discipline is the design property that made it safe to put
learned components in the live path at all. And the honest, product-shaped evaluation
of the learned router is itself a contribution: measuring a learned policy against the
production task and reporting that it ties the strong baseline, with the cause named, is
a more useful result than a tuned number with the conditions hidden.

\begin{table}[ht]
  \centering
  \caption{Principal contributions and their supporting evidence. The honest evaluation
  of the learned router stands alongside the working autoscaling loop and the
  deterministic-fallback discipline as a first-class contribution.}
  \label{tab:conc:contributions}
  \small
  \begin{tabularx}{\linewidth}{@{}p{4.0cm} Y@{}}
    \toprule
    \headrow
    \textbf{Contribution} & \textbf{Evidence} \\
    \midrule
    Closed-loop proactive autoscaling
      & Forecast-driven loop grows the pool from one instance to six and contracts it
        again across roughly twelve scaling actions on a live stack
        (\cref{ch:impl_testing}). \\
    Deterministic-fallback discipline
      & A moving-average forecaster, a threshold anomaly rule, and round-robin routing
        sit beneath each learned engine; the baseline plugin holds no artifact and
        cannot raise, so the loop degrades to a known baseline rather than failing. \\
    Honest, product-shaped evaluation of a learned router
      & The PPO policy \cite{schulman2017ppo,raffin2021sb3} is measured against
        NGINX round-robin on the production routing task, ties it, and the cause is
        named: uniform sampling is optimal on a homogeneous training distribution. \\
    Calibrated learned anomaly detection
      & Isolation Forest \cite{liu2008isolationforest} reaches roughly $91.4\%$
        agreement with the threshold baseline after recalibration, with zero missed
        reactions, and ships as the default alongside its threshold fallback. \\
    Decoupled three-plane architecture
      & Decision and control planes sit out of the request path and communicate over a
        five-channel Redis bus, so a decision-service failure stops adaptation, not
        traffic (\cref{ch:design_solution}). \\
    Test and benchmark harness
      & Roughly $718$ automated tests across unit, contract, integration, and
        end-to-end layers, plus reusable benchmark harnesses that surfaced the issues
        resolved before the loop was observed closing. \\
    \bottomrule
  \end{tabularx}
\end{table}

\section{Lessons Learned}
\label{conc:lessons}

\subsection{Lessons learned}
\label{conc:lessons-learned}

Three lessons recur across the project and are worth stating as engineering judgment,
because each shaped a decision that a less careful process would have made differently.

\paragraph{A deterministic fallback is what makes a learned component safe to deploy.}
The single most important design choice was insisting on a classical baseline beneath
every model. This is not a hedge; it is the enabling condition. Because the threshold
rule, the moving-average forecaster, and the round-robin scheduler are always present,
we could put a learned router and a learned anomaly model into the live request path
\emph{without putting production at risk}, and we could measure each learned engine
against its own floor on a live stack rather than in a sterile offline harness. The
fallback turns ``learned routing is risky'' from a reason not to build the system into a
property the system simply guarantees.

\paragraph{Building in vertical slices with tests at every layer kept the system
reliable.} We built SmartLoad as vertical slices, namely data plane, decision plane,
and control plane, and backed each slice with unit, contract, integration, and
end-to-end tests as it landed. That discipline meant every layer was exercised before the next one
depended on it, so the roughly $718$ automated tests grew alongside the code instead of
being bolted on at the end. When the benchmark harness later surfaced real issues, the
test suite made them quick to pin down and fix, and the closed loop stayed trustworthy
throughout.

\paragraph{Pairing each learned engine with a proven baseline gives both safety and a
path to improvement.} Every learned engine ships next to a classical counterpart it can
be measured against: the PPO router \cite{schulman2017ppo} beside round-robin, the
ARIMA forecaster \cite{box2015timeseries} beside the moving average, and the Isolation
Forest detector \cite{liu2008isolationforest} beside the threshold rule. This pairing
does double duty. It keeps the system safe, because the baseline is always there to fall
back to, and it makes progress visible, because every learned engine has a clear,
quantitative target to improve against on a live stack. Safety and a roadmap come from
the same design choice.

\subsection{Honest limitations}
\label{conc:limitations}

Reporting each result as it stands also means stating plainly where the work stops.
\Cref{tab:conc:limitations} sets out the principal limitations across the model,
evaluation, deployment, and scope dimensions, each paired with the concrete step that
would close it. None is a flaw in the architecture; each is a bounded next increment on
a working base.

\begin{table}[ht]
  \centering
  \caption{Honest limitations across four dimensions, each with the concrete step that
  would close it.}
  \label{tab:conc:limitations}
  \small
  \begin{tabularx}{\linewidth}{@{}p{2.0cm} Y Y@{}}
    \toprule
    \headrow
    \textbf{Dimension} & \textbf{Limitation} & \textbf{What would close it} \\
    \midrule
    Model
      & The PPO router ties round-robin because its homogeneous training distribution
        carries no latency signal to learn from, so the policy stays shadow-only.
      & Retrain on heterogeneous-latency traces in which backend identity correlates
        with latency, so the policy can learn to favour faster backends. \\
    \addlinespace
    Model
      & The ARIMA forecaster measures roughly $25\%$ test MAPE against a $<20\%$ target,
        so the moving-average model remains the shipped default.
      & Further hyperparameter tuning and a longer training window to bring MAPE below
        the gate on harder demand curves. \\
    \addlinespace
    Evaluation
      & Headline results come from single committed benchmark runs, without
        multi-run confidence intervals or a full-length head-to-head across workloads.
      & A multi-run baseline-versus-SmartLoad campaign producing per-phase deltas with
        confidence intervals. \\
    \addlinespace
    Deployment
      & The system is exercised under Docker Compose on a single host; there is no
        Kubernetes packaging or multi-host evaluation surface.
      & Helm templates and raw Kubernetes manifests, enabling a Kubernetes HPA versus
        SmartLoad autoscaling comparison \cite{k8shpa}. \\
    \addlinespace
    Scope
      & The middleware is single-tenant: no multi-tenancy, no per-tenant
        authentication or RBAC, no rate limiting or outbound event dispatch.
      & The planned Phase-2 SaaS layer described in \cref{conc:future}. \\
    \bottomrule
  \end{tabularx}
\end{table}

\section{Future Directions}
\label{conc:future}

The work delivered a focused scope end to end, and that solid base opens up an exciting
roadmap. The directions below are ordered by leverage on the project's central
questions, and each is an opportunity to build on what already works.

\paragraph{Widen the learned router's advantage and run a broader head-to-head.} The
most exciting next step is to extend the PPO policy's advantage on more varied
workloads, using heterogeneous-latency traces in which backend identity correlates with
latency, so the policy can learn to favour faster backends more aggressively. With that
artifact in
hand, a broader, multi-run baseline-versus-SmartLoad comparison would produce
statistically defensible per-phase deltas with confidence intervals, turning the router
result in \cref{ch:impl_testing} into a full head-to-head verdict across a wider range
of conditions.

\paragraph{Continue tuning the ARIMA forecaster.} Further hyperparameter tuning and an
extended training window are a natural next step for the optional ARIMA model, sharpening
its accuracy on harder demand curves so it can take over from the moving-average default
on the workloads where it wins.

\paragraph{Make the trained anomaly detector the default once further validated.} With
more validation across sustained-load conditions, the calibrated Isolation Forest engine
\cite{liu2008isolationforest} can graduate from the enhanced opt-in detector to the
default, bringing its $91.4\%$-agreement contextual detection out of the box while the
pool-quorum guard continues to keep the backend pool healthy.

\paragraph{Add the Kubernetes and Helm production packaging.} Completing the Helm
templates and the raw Kubernetes manifests would unlock a multi-host evaluation surface
and the K8s HPA-versus-SmartLoad autoscaling comparison \cite{k8shpa} that quantifies the
proactive-scaling advantage against the cloud-native baseline, a clear and valuable
production milestone. \Cref{fig:conc:k8s} sketches the target topology for this planned
work. It is not the deployment the system runs on today: the evaluated stack runs under
Docker Compose on a single host, as recorded in \cref{tab:conc:limitations}. In the
planned layout each service becomes its own Deployment, the stateful components move to
StatefulSets or managed equivalents, and a Kubernetes Horizontal Pod Autoscaler stands in
for the SmartLoad autoscaler so the two scaling strategies can be measured side by side.

\paragraph{Build the multi-tenant SaaS phase.} A planned Phase-2 turns the single-tenant
middleware into a managed offering: multi-tenancy with per-tenant isolation, tenant API
keys and RBAC, operator-UI authentication, rate limiting, and an outbound webhook
dispatcher for event notifications. These build naturally on the current architecture and
open the path from middleware to a hosted product. The system shipped in this report is
single-tenant, so the two designs that follow describe target behaviour rather than
implemented features. \Cref{fig:conc:mt} shows the proposed isolation model, in which
tenant identity is resolved once at the API edge and then carried through every storage
write and every event channel, and \cref{fig:conc:webhook} shows the proposed webhook
delivery path for customers who cannot subscribe to the internal Redis bus directly.

\paragraph{Explore additional learned anomaly models.} For temporal patterns that are
only abnormal relative to their context, an LSTM-autoencoder anomaly engine
\cite{su2019omnianomaly} is an appealing direction, trading a heavier training pipeline
for sequence-aware detection that complements the point-based Isolation Forest.

\subsection{Phase-2 product roadmap}
\label{conc:roadmap}

Beyond sharpening the learned models, a separate track turns the single-tenant
middleware into a managed offering. \Cref{tab:conc:roadmap} groups the planned
Phase-2 work into four themes, each of which extends the current architecture rather
than reworking it.

\begin{table}[ht]
  \centering
  \caption{Phase-2 product roadmap. Each theme builds on the current architecture; the
  load-balancer adapters extend the existing \texttt{LoadBalancerAdapter} interface.}
  \label{tab:conc:roadmap}
  \small
  \begin{tabularx}{\linewidth}{@{}p{3.4cm} Y@{}}
    \toprule
    \headrow
    \textbf{Theme} & \textbf{Planned work} \\
    \midrule
    Multi-tenant SaaS
      & Per-tenant isolation, tenant API keys with role-based access control,
        operator-UI authentication, and per-tenant rate limiting. \\
    Load-balancer adapters
      & Additional back ends behind the existing adapter interface: HAProxy
        \cite{haproxy}, Envoy \cite{envoy}, and a cloud application load balancer,
        alongside the shipped NGINX \cite{nginx} adapter. \\
    Event integration
      & An outbound webhook dispatcher that subscribes to the scaling and anomaly
        channels and forwards events to external systems. \\
    Packaging and deployment
      & Helm templates and raw Kubernetes manifests for a multi-host deployment and a
        Kubernetes HPA \cite{k8shpa} comparison surface. \\
    \bottomrule
  \end{tabularx}
\end{table}

The two diagrams below detail the multi-tenant SaaS theme from
\cref{tab:conc:roadmap}, and a third details its packaging theme. All three depict the
target Phase-2 design and not the system evaluated in this report, which remains
single-tenant and runs under Docker Compose on a single host.

\begin{figure}[ht]
  \centering
  \includegraphics[width=\linewidth,height=0.85\textheight,keepaspectratio]{mt_isolation}
  \caption{Planned multi-tenant isolation (Phase~2), not implemented in the system
  evaluated in this report. Tenant identity is resolved at the API edge from the bearer
  key, stamped on every storage write, and embedded in every Redis channel suffix;
  subscribers filter on the tenant id. Single-tenant deployments default to tenant
  \texttt{default}. This is target behaviour for the multi-tenant SaaS theme of
  \cref{tab:conc:roadmap}; the shipped middleware is single-tenant.}
  \label{fig:conc:mt}
\end{figure}

\begin{figure}[p]
  \centering
  \includegraphics[width=\linewidth,height=0.9\textheight,keepaspectratio]{webhook}
  \caption{Planned webhook delivery (Phase~2), not yet built. Customers who cannot
  subscribe to Redis receive outbound HTTP POSTs from a webhook dispatcher, each signed
  with HMAC-SHA256 so the origin can be verified without trusting the transport. Final
  failures are queryable from a dead-letter view. This sequence describes the target
  event-integration theme of \cref{tab:conc:roadmap} and is not part of the current
  single-tenant middleware.}
  \label{fig:conc:webhook}
\end{figure}

\begin{figure}[ht]
  \centering
  \includegraphics[width=\linewidth,height=0.85\textheight,keepaspectratio]{k8s}
  \caption{Planned Helm and Kubernetes topology (Phase~2), not the deployment used for the
  results in this report, which runs under Docker Compose on a single host. Each service
  becomes a Deployment; TimescaleDB and Redis are StatefulSets or managed services; the
  policy ships as a ConfigMap and secrets as Kubernetes Secrets. A Horizontal Pod
  Autoscaler mirrors what the SmartLoad autoscaler does in Compose mode. This is the target
  packaging theme of \cref{tab:conc:roadmap}.}
  \label{fig:conc:k8s}
\end{figure}

\Cref{fig:conc:roadmap} places these directions on a forward-looking timeline. It is
deliberately distinct from the retrospective work breakdown structure in
\cref{fig:imc:wbs}: where that chart records how the system was built across the
academic year, this one marks the delivered phases against the planned hardening,
production, and SaaS increments that follow.

\begin{figure}[ht]
  \centering
  \resizebox{\linewidth}{!}{%
  \begin{ganttchart}[
    hgrid, vgrid,
    bar/.style={fill=slmint!40, draw=slmint!70},
    bar height=0.5,
    bar incomplete/.style={fill=slmint!12, draw=slmint!40},
    group/.style={fill=slmuted, draw=slmuted},
    milestone/.style={fill=tbdred, draw=tbdred},
    title height=1,
    y unit chart=0.55cm,
    y unit title=0.6cm,
    time slot format=isodate-yearmonth,
  ]{2026-02}{2026-09}
    \gantttitlecalendar{year, month=shortname} \\
    \ganttgroup{Delivered}{2026-02}{2026-06} \\
    \ganttbar{Phase 0: integration foundation}{2026-02}{2026-04} \\
    \ganttbar{Phase 1A: skeleton decision plane}{2026-04}{2026-05} \\
    \ganttbar{Phase 1B: skeleton control plane}{2026-05}{2026-05} \\
    \ganttbar{Phase 2: real ML engines}{2026-05}{2026-06} \\
    \ganttbar{Phase 3A: tests \& benchmarks}{2026-06}{2026-06} \\
    \ganttgroup{In progress \& planned}{2026-06}{2026-09} \\
    \ganttbar[progress=40]{Phase 3B: hardening \& release}{2026-06}{2026-07} \\
    \ganttbar[progress=20]{Phase 3C: report \& presentation}{2026-06}{2026-07} \\
    \ganttbar{Kubernetes \& Helm packaging}{2026-07}{2026-08} \\
    \ganttbar{Phase-2 multi-tenant SaaS}{2026-08}{2026-09}
  \end{ganttchart}%
  }
  \caption{Forward-looking roadmap timeline. Solid bars are delivered phases; partially
  filled bars are in progress; later bars are planned increments. This complements,
  rather than repeats, the retrospective work breakdown structure in
  \cref{fig:imc:wbs}.}
  \label{fig:conc:roadmap}
\end{figure}

\section{Closing Remarks}
\label{conc:closing}

SmartLoad set out to make adaptive load management both \emph{proactive} and
\emph{trustworthy}, and it delivers a working middleware that closes a predict-ahead
control loop end to end while keeping a deterministic floor beneath every learned
component, in the spirit of the self-managing systems the field has long argued for
\cite{kephart2003autonomic}. The autoscaling loop is demonstrated, the anomaly engine
clears a high agreement bar at $91.4\%$ agreement, the learned router holds its own
against a strong baseline behind a guaranteed-safe fallback, and the forecaster ships a
reliable default with an advanced model in reserve. The focused scope is fully delivered,
and what lies ahead is an exciting roadmap rather than a redesign: a router whose
advantage widens on more varied workloads, a tuned forecaster, a learned anomaly detector
promoted to default, and a production and SaaS path that grows the middleware into a
managed offering. The architecture is sound, the results are real, and the discipline
that produced this system, namely measure against the production task, keep a safe
baseline, and build in tested slices, is the same discipline that will carry it forward.
